\section{Ferramentas e Validação de \textit{Software}}
\label{sec:ferramentas_software}

Por período considerável, o \textit{firmware} e a aplicação de supervisão
precisaram ser escritos e depurados antes de a placa de medição estar
fabricada e validadas em bancada. Esta seção resume as ferramentas adotadas
para viabilizar esse desenvolvimento desacoplado do \textit{hardware} físico.

Os esquemáticos e o roteamento das placas de circuito impresso foram
desenvolvidos no \textit{KiCad} (\textit{software} de CAD eletrônico de código
aberto); o \textit{front-end} analógico de medição de corrente e tensão de
fase (\autoref{sec:same_tensao}) foi previamente simulado no
\textit{LTspice} antes da fabricação, reduzindo o número de iterações físicas
em placas de fenolite (\autoref{sec:padronizacao_placas}).

Para o \textit{firmware}, o \textit{Zephyr RTOS} oferece um alvo de compilação
alternativo chamado \textit{native\_sim} (\textit{simulação nativa}): em vez de
gerar um binário para o núcleo RISC-V da placa, o mesmo código de aplicação é
compilado diretamente para a arquitetura do computador de desenvolvimento e
roda como um processo comum do Linux, sem emulação de processador. A rede
permanece real (o processo abre \textit{sockets} UDP de verdade, em
\textit{loopback}), e os periféricos físicos ausentes (ADC, PWM do ESC,
contadores de RPM) são substituídos por geradores de valor sintético dentro do
mesmo arquivo de \textit{hardware} usado na placa real, sob diretiva de
compilação condicional, nunca em uma cópia paralela. Essa prática permitiu
encontrar, ainda em desenvolvimento, defeitos de protocolo que só apareceriam
na integração física: um exemplo concreto foi a divergência do tamanho da
estrutura de comando (\texttt{CMD}) entre o compilador do microcontrolador
(\textit{enums} de 1 \textit{byte}) e o do computador hospedeiro (\textit{enums}
de 4 \textit{bytes} por padrão), detectada justamente ao validar contra o
\textit{native\_sim}. A limitação da técnica é não exercitar nada que dependa
do comportamento físico do \textit{hardware} (temporização de periférico, DMA,
pinagem); ausência de defeito nesse modo não substitui a validação final em
bancada.

A mesma restrição levou a aplicação de supervisão a suportar três formas
alternativas de conexão com cada placa, escolhidas em tempo de execução sem
alterar o restante do código da interface, conforme a
\autoref{fig:modos_conexao} e a \autoref{tab:modos_conexao}.

\begin{figure}[H]
    \centering
    \caption{Transições entre os modos de conexão da aplicação de supervisão}
    \label{fig:modos_conexao}
    \begin{tikzpicture}[
        font=\small, >={Latex},
        estado/.style={draw, thick, rounded corners=6pt, align=center,
            minimum height=1.25cm, text width=4.8cm, drop shadow={opacity=0.15}},
        rot/.style={font=\scriptsize, fill=white, inner sep=1.5pt}
    ]
        \node[estado, fill=gray!12, text width=2.8cm, minimum height=2.6cm] (sem)
            {\textbf{Desconectado}\\{\footnotesize nenhum \textit{link} ativo}};
        \node[estado, fill=green!15, draw=green!45!black] (hw) at (7.6,1.9)
            {\textbf{\textit{Hardware} real}\\{\footnotesize placa física, via Ethernet}};
        \node[estado, fill=blue!12, draw=blue!50!black] (mock) at (7.6,0)
            {\textbf{Simulador (\textit{mock})}\\{\footnotesize processo no próprio PC}};
        \node[estado, fill=orange!18, draw=orange!70!black] (nat) at (7.6,-1.9)
            {\textbf{\textit{native\_sim}}\\{\footnotesize \textit{firmware} compilado para o PC}};

        \draw[<->, thick, green!45!black]  ([yshift=0.8cm]sem.east)  -- node[rot, sloped] {conectar / desconectar} (hw.west);
        \draw[<->, thick, blue!50!black]   (sem.east)                 -- node[rot] {conectar / desconectar} (mock.west);
        \draw[<->, thick, orange!70!black] ([yshift=-0.8cm]sem.east) -- node[rot, sloped] {conectar / desconectar} (nat.west);
    \end{tikzpicture}

    Fonte: Autor (2026)
\end{figure}

\begin{table}[h]
\centering
\caption{Modos de conexão da aplicação de supervisão com as placas.}
\label{tab:modos_conexao}
\begin{tabular}{p{2.2cm}p{4.4cm}p{5.6cm}}
\hline
\textbf{Modo} & \textbf{Onde executa} & \textbf{Uso metodológico} \\
\hline
\textit{Hardware} real & Placa física, via Ethernet & Validação final e ensaios de bancada \\
Simulador Python (\textit{mock}) & \textit{Thread} que finge o protocolo por correspondência de nomes & Ajuste de interface sem placa disponível; não valida o protocolo real \\
\textit{native\_sim} & \textit{Firmware} real, compilado para POSIX, em \textit{loopback} & Validação do protocolo e da lógica do \textit{firmware} sem bancada \\
\hline
\end{tabular}
\end{table}

Por fim, todo o código do trabalho (projetos de \textit{hardware} (KiCad),
\textit{firmware} (C/Zephyr) e aplicação de supervisão em Python) é mantido
em um único repositório Git, o que simplifica a rastreabilidade de mudanças
que afetam simultaneamente os dois lados do protocolo de rede
(\autoref{fig:yaml_example_config}).
